\section{Relationship to Mergeable Summaries}\label{app:mergeable}

This appendix formalizes the connection between our Oracle-Preserving Summarization (OPS) framework and the theory of \emph{mergeable summaries} developed for streaming and distributed databases \citep{Agarwal2013MergeableTODS}. We show that OPS strictly generalizes mergeable summaries: when the oracle $\fstar$ is symmetric and the summarizer $g$ is deterministic, our Parentwise Compatibility law~\eqref{eq:leaf_compatibility} is identical to the mergeability condition. When $\fstar$ is order-sensitive or $g$ is stochastic, OPS extends the notion of mergeability to the free monoid of strings by replacing hand-designed merge operators with learned semantic merges implemented by an LLM.

\subsection{Formal Setup of Mergeable Summaries}

In the classical setting, data arrives as a collection of items, and the objective is to maintain a compact synopsis that can be merged efficiently. Let $\mathcal{D}$ denote datasets (typically multisets) and $\mathcal{S}$ the sketch space with $|\mathcal{S}|\ll|\mathcal{D}|$.
\begin{itemize}
    \item \textbf{Data:} $D_1, D_2 \in \mathcal{D}$.
    \item \textbf{Operation:} Multiset union $\uplus$, which is associative and commutative.
    \item \textbf{Summary:} A function $S:\mathcal{D}\to\mathcal{S}$.
    \item \textbf{Query:} $Q$ maps a sketch to the target statistic, with distortion bounded by $\epsilon$.
\end{itemize}

\begin{definition}[Mergeable Summary]
A summary $S$ is mergeable if there exists $\mathcal{M}:\mathcal{S}\times\mathcal{S}\to\mathcal{S}$ such that for all $D_1,D_2$,
\[
    d_{\Y}\!\left(Q\!\left(\mathcal{M}(S(D_1),S(D_2))\right),\, Q(D_1\uplus D_2)\right) \le \epsilon.
\]
\end{definition}

The guarantee holds uniformly over arbitrary merge trees, allowing sketches to be composed without revisiting the raw inputs.

\subsection{Mapping OPS to Mergeable Summaries}

Table~\ref{tab:mergeable_mapping} aligns the notation used in Section~\ref{sec:mechanism} with the streaming literature. The distinction is that OPS works over the free monoid of strings $(\Strings,\concat)$, where $\concat$ is generally non-commutative, and uses an LLM summarizer $g$ instead of a fixed algebraic operator.

\begin{table}[t]
    \centering
    \begin{tabular}{p{0.23\linewidth}p{0.33\linewidth}p{0.33\linewidth}}
        \toprule
        \textbf{Component} & \textbf{OPS (this paper)} & \textbf{Mergeable summaries} \\
        \midrule
        Input & Strings $u,v\in\Strings$ & Datasets $D_1,D_2\in\mathcal{D}$ \\
        Combination & Concatenation $\concat$ (non-commutative) & Union $\uplus$ (commutative) \\
        Summarizer & $g:\Strings\rightsquigarrow\Strings$ & $S:\mathcal{D}\to\mathcal{S}$ \\
        Task/query & Oracle $\fstar:\Strings\to\Y$ & Query $Q$ \\
        Merge algorithm & $g(g(u)\concat g(v))$ & $\mathcal{M}(S(D_1),S(D_2))$ \\
        Consistency law & Parentwise compatibility \eqref{eq:leaf_compatibility} & Mergeability condition \\
        \bottomrule
    \end{tabular}
    \caption{Mapping OPS to the mergeable summaries framework.}
    \label{tab:mergeable_mapping}
\end{table}

\subsection{Equivalence in the Commutative Case}

Suppose the task ignores order, such as counting keyword occurrences or extracting the set of unique actors. In this case, $\fstar$ acts as a homomorphism from $(\Strings,\concat)$ to a commutative monoid $(\Y,+)$:
\[
    \fstar(u\concat v) = \fstar(u) + \fstar(v),
\]
up to the oracle metric $\metric$.

Parentwise compatibility then states that re-summarizing the child summaries is inert at the oracle:
\[
    \E\!\Big[\metric\!\Big(\fstar\big(g(g(u)\concat g(v))\big),\, \fstar(g(u)\concat g(v))\Big)\Big]=0.
\]
When $g$ is deterministic, the expectation drops and the condition becomes pointwise.

\begin{proposition}[Reduction to Mergeable Sketches]
\label{prop:mergeable_reduction}
Let $\fstar$ be commutative so that $\fstar(u\concat v)=\fstar(v\concat u)$, and let $g$ be deterministic. If $g$ satisfies exact Parentwise Compatibility~\eqref{eq:leaf_compatibility} on every realized merge, then $g$ is a mergeable summary for $\fstar$ in the sense of \citet{Agarwal2013MergeableTODS}.
\end{proposition}

\begin{proof}
Define the merge algorithm $\mathcal{M}(s_1,s_2)=g(s_1\concat s_2)$, where $s_i=g(u_i)$ are summaries. For any $u,v$ we compare two routes: (i) concatenate the summaries $s_u,s_v$ and apply $\fstar$, yielding $\fstar(g(u)\concat g(v))$; (ii) merge the summaries via $\mathcal{M}$ and apply $\fstar$, yielding $\fstar(g(g(u)\concat g(v)))$. Parentwise compatibility enforces equality between these two values, so the mergeability condition holds with $\epsilon=0$. When the oracle is commutative and Condition~\eqref{eq:leaf_sufficiency} holds on the leaves, the concatenated summaries themselves agree with the concatenated raw text under $\fstar$, recovering the classical guarantee.
\end{proof}

Note that, unlike classical sketches where data arrive already as sketch inputs, Condition~(C2) in our framework also constrains the \emph{raw} concatenations $u\concat v$ to agree with their summarized counterparts. The parsing step from text to summary is therefore part of the algebra itself, which makes the OPS guarantees strictly stronger than standard mergeability claims.

\begin{remark}[The universal merge]\label{rem:universal_merge}
Classical sketches design bespoke $\mathcal{M}$ operators (vector addition for Count-Min, rank pruning for quantile sketches). OPS instead uses a universal merge: concatenate the summaries and apply the same summarizer again. The summarizer learns the algebra needed by the downstream oracle during its initial reduction pass, so no task-specific merge code is required.
\end{remark}

\subsection{Generalization to Non-Commutative Tasks}

The OPS framework extends mergeability to the free monoid where $u\concat v\neq v\concat u$. Narrative queries---``Did the protest in paragraph $A$ trigger the crackdown in paragraph $B$?''---depend on order and context. Standard sketches over multisets cannot express these dependencies because the union operator discards order. Parentwise compatibility enforces a non-commutative analogue: the interaction between summarized blocks mimics the interaction between the original texts when evaluated by $\fstar$. This is precisely the condition audited in Figure~\ref{fig:local_compatibility}.

\subsection{Error Propagation and Bounds}

Mergeable summaries typically admit error bounds that grow logarithmically with tree height. Theorems~\ref{thm:one-pass} and~\ref{thm:multi-round} give an analogous ``zero-error'' propagation: if local merges satisfy Parentwise Compatibility exactly (in expectation), the global error collapses to zero regardless of tree shape. When local failures are bounded by $\epsilon$, our deployment bound~\eqref{eq:error_budget} plays the role of classical error accumulation: the total distortion is controlled by a weighted sum of leaf, merge, and stabilization failure rates. This is the semantic counterpart to bounding sketch errors with respect to the number of merge operations; it justifies the probabilistic audit in Section~\ref{sec:audit}, where random sampling over the summary tree produces a high-confidence certificate for the entire hierarchy.
